\section{Resultados}\label{resultados}

\subsection{Pregunta de
investigación}\label{pregunta-de-investigaciuxf3n}

¿Qué tan efectivo es Naive Bayes en clasificación de texto?

\subsection{Cadena de búsqueda}\label{cadena-de-buxfasqueda}


\begin{verbatim}
("naive bayes" OR "naive bayes classifier" OR "naive bayes classification" OR
"bayes classifier" OR "multinomial naive bayes" OR "bernoulli naive bayes")
AND
("text classification" OR "document classification" OR "text categor*" OR
"text mining" OR "text analytics" OR "natural language processing" OR "NLP" OR
"sentiment analysis" OR "topic classification" OR "document categor*")
\end{verbatim}


\subsection{Clasificación de los 44
artículos}\label{clasificaciuxf3n-de-los-44-artuxedculos}

A continuación se presenta la clasificación de los 44 artículos identificados en las fuentes, evaluados según su relevancia para investigar la efectividad de \textbf{Naive Bayes (NB)} en la clasificación de textos.

Los criterios utilizados son:
\begin{itemize}
    \item \textbf{Alto:} El artículo se centra específicamente en mejorar el algoritmo, comparar su rendimiento de forma exhaustiva o proponer variantes fundamentales.
    \item \textbf{Medio:} El artículo aplica NB a un dominio, idioma o tarea específica (por ejemplo, medicina, árabe, detección de clickbait), priorizando la aplicación sobre el estudio teórico de la efectividad general del algoritmo.
\end{itemize}

\begin{longtable}{|c|c|p{9cm}|}
\caption{Clasificación de artículos según su relevancia para el estudio de Naive Bayes en clasificación de textos}
\label{tab:clasificacion-articulos} \\
\hline
\textbf{\#} & \textbf{Nivel de Relevancia} & \textbf{DOI o URL del Artículo} \\
\hline
\endfirsthead

\hline
\textbf{\#} & \textbf{Nivel de Relevancia} & \textbf{DOI o URL del Artículo} \\
\hline
\endhead

[1]  & \textbf{Alto}  & \url{10.1109/TKDE.2006.180} \\
\hline
[2]  & \textbf{Medio} & \url{10.1109/SNPD.2019.8935711} \\
\hline
[3]  & \textbf{Alto}  & \url{10.1007/11871637_49} \\
\hline
[4]  & \textbf{Alto}  & \url{10.1145/2911451.2911547} \\
\hline
[5]  & \textbf{Alto}  & \url{10.1016/J.ENGAPPAI.2016.02.002} \\
\hline
[6]  & \textbf{Alto}  & \url{https://doi.org/10.1007/3-540-45683-X_45} \\
\hline
[7]  & \textbf{Alto}  & \url{https://doi.org/10.1007/978-3-540-30586-6_76} \\
\hline
[8]  & \textbf{Alto}  & \url{https://doi.org/10.1007/978-3-540-71496-5_73} \\
\hline
[9]  & \textbf{Alto}  & \url{https://doi.org/10.1007/978-3-540-30549-1_43} \\
\hline
[10] & \textbf{Medio} & \url{https://doi.org/10.1007/11735106_17} \\
\hline
[11] & \textbf{Alto}  & \url{10.1016/j.eswa.2008.06.054} \\
\hline
[12] & \textbf{Alto}  & \url{10.1016/j.proeng.2011.08.404} \\
\hline
[13] & \textbf{Alto}  & \url{10.1016/j.patrec.2015.07.028} \\
\hline
[14] & \textbf{Alto}  & \url{10.1016/j.patrec.2008.11.013} \\
\hline
[15] & \textbf{Alto}  & \url{10.1016/j.knosys.2016.02.017} \\
\hline
[16] & \textbf{Alto}  & \url{10.1109/ISIP.2008.54} \\
\hline
[17] & \textbf{Alto}  & \url{10.1016/j.asoc.2020.106652} \\
\hline
[18] & \textbf{Medio} & \url{10.1016/j.knosys.2010.04.004} \\
\hline
[19] & \textbf{Alto}  & \url{10.1016/j.neucom.2018.07.002} \\
\hline
[20] & \textbf{Alto}  & \url{https://dl.acm.org/doi/10.5555/1619645.1619732} \\
\hline
[21] & \textbf{Medio} & \url{10.32604/cmes.2023.026812} \\
\hline
[22] & \textbf{Medio} & \url{10.55151/ijeedu.v5i1.99} \\
\hline
[23] & \textbf{Medio} & \url{10.62441/nano-ntp.v20iS1.77} \\
\hline
[24] & \textbf{Medio} & \url{10.38094/jastt61237} \\
\hline
[25] & \textbf{Medio} & \url{10.3390/bdcc7030127} \\
\hline
[26] & \textbf{Medio} & \url{10.1063/5.0110821} \\
\hline
[27] & \textbf{Medio} & \url{10.1016/j.fraope.2025.100373} \\
\hline
[28] & \textbf{Alto}  & \url{10.1007/s11634-023-00536-8} \\
\hline
[29] & \textbf{Medio} & \url{10.1051/e3sconf/202346502048} \\
\hline
[30] & \textbf{Medio} & \url{10.1109/ACCESS.2024.3487311} \\
\hline
[31] & \textbf{Alto}  & \url{10.26615/978-954-452-056-4_151} \\
\hline
[32] & \textbf{Medio} & \url{10.1088/1742-6596/1641/1/012043} \\
\hline
[33] & \textbf{Alto}  & \url{10.5220/0010890400003122} \\
\hline
[34] & \textbf{Alto}  & \url{10.1109/TKDE.2016.2522427} \\
\hline
[35] & \textbf{Alto}  & \url{10.3390/math9192378} \\
\hline
[36] & \textbf{Alto}  & \url{10.1109/TKDE.2010.160} \\
\hline
[37] & \textbf{Alto}  & \url{10.1007/978-3-642-41278-3_24} \\
\hline
[38] & \textbf{Alto}  & \url{10.1109/ACCESS.2020.2968984} \\
\hline
[39] & \textbf{Alto}  & \url{10.1109/TKDE.2016.2563436} \\
\hline
[40] & \textbf{Medio} & \url{10.1093/bioinformatics/btl350} \\
\hline
[41] & \textbf{Medio} & \url{10.4108/eai.12-10-2019.2296530} \\
\hline
[42] & \textbf{Medio} & \url{10.35940/ijeat.A1610.109119} \\
\hline
[43] & \textbf{Medio} & \url{10.1186/1471-2105-8-269} \\
\hline
[44] & \textbf{Alto}  & \url{10.4067/S0718-09342011000300004} \\
\hline

\end{longtable}

\subsection{Gerarquización (Alto, Medio, Bajo) de los 44 artículos en Notebook}

\begin{figure}[H]
    \centering
    
\includegraphics[width=0.95\textwidth]{\detokenize{(trozos) (6.2) Capturas de pantalla (NotebookLM).png/trozo_0.png}}
    \caption{Notebook 1}
    \label{fig:notebook1}
\end{figure}

\begin{figure}[H]
    \centering
    \includegraphics[width=0.95\textwidth]{\detokenize{(trozos) (6.2) Capturas de pantalla (NotebookLM).png/trozo_1.png}}
    \caption{Notebook 2}
    \label{fig:notebook2}
\end{figure}

\begin{figure}[H]
    \centering
    \includegraphics[width=0.95\textwidth]{\detokenize{(trozos) (6.2) Capturas de pantalla (NotebookLM).png/trozo_2.png}}
    \caption{Notebook 3}
    \label{fig:notebook3}
\end{figure}

\subsection{Los 44 artículos importados en Mendeley}

\begin{figure}[H]
    \centering
    \includegraphics[width=0.95\textwidth]{\detokenize{(6.3) Captura en Mendeley.png}}
    \caption{Los 44 artículos importados en Mendeley}
    \label{fig:mendeley44}
\end{figure}